\section{Methodology}
\label{sec:methodology}

\subsection{Shared Routing Model}

The operating area is an $N\times N$ grid with coordinates $(row,col)$.
Movement is eight-connected. Orthogonal edges cost $1$ and diagonal edges cost
$\sqrt{2}$; diagonal corner cutting is allowed. A destination cell must be
inside the grid and traversable. Optional non-negative traversal penalties
represent wind, energy demand, or penalized no-fly regions. Hard no-fly cells
are removed from the traversable graph. All planners obtain successors and
edge costs from the same environment function.

An obstacle may toggle on a fixed period, toggle stochastically under a seeded
process, or move along a cyclic path. Every method follows the same
\texttt{post\_move\_observed} information contract: it selects and executes a
move from the current state, dynamics then advance, and the changed state is
available before the next decision. Terminal moves produce no unobservable
event, and methods receive no future schedule. If an obstacle appears in the
occupied cell, the benchmark permits the UAV to exit but not re-enter that
cell. This legacy contract is shared across methods and reported as a
limitation.

\subsection{Classical Methods}

\textbf{Dijkstra} performs a complete shortest-path search initially and after
each observed graph change. It is an interpretable full-replanning reference
with complexity $O((V+E)\log V)$ under a binary heap.

\textbf{A*} uses octile distance, the admissible heuristic for the shared
eight-connected movement geometry \cite{hart1968astar}. It shares Dijkstra's
runner and replanning triggers, isolating heuristic pruning as the only
algorithmic change.

\textbf{D* Lite} maintains $g$ and one-step-lookahead $rhs$ values and updates
only vertices affected by changed edges \cite{koenig2005dstarlite}. It receives
the same change notifications as the full replanners. Its cumulative expansion
count includes the initial computation and all incremental repairs.

\subsection{Reinforcement Learning Policy}
\label{sec:rl-policy}

The learned method is Double DQN \cite{vanhasselt2016double} with HER
\cite{andrychowicz2017her}. The policy uses a two-layer MLP with 128 units per
layer. Its default 61-dimensional observation concatenates normalized UAV
position and goal displacement (4), a flattened $7\times7$ local occupancy
window (49), and a one-hot previous action (8). The full-observation ablation
replaces the local window with the complete grid.

The sparse reward is $+1$ at the goal, $-1$ on collision, and $-0.1$ otherwise.
Potential shaping uses
\begin{equation}
F(s,s')=\gamma\Phi(s')-\Phi(s), \qquad
\Phi(s)=-d_{\mathrm{octile}}(s,g)/d_{\max},
\end{equation}
with $\gamma=0.99$ \cite{ng1999shaping}. Octile distance is deliberately
independent of obstacle state. This is necessary for valid HER relabeling:
recomputing a replayed reward from the environment's obstacle state at sampling
time would not reproduce the state under which the transition was collected.
Shortest-path distance is retained for evaluation metrics, not used as the
dynamic HER potential.

All learned variants use a learning rate of $10^{-3}$, replay capacity
$100{,}000$, 5,000 learning-start steps, batch size 256, discount $0.99$, one
gradient update every two environment steps, and a target-network update every
1,000 steps. Exploration decays linearly from $1.0$ to $0.05$ over the first
half of training. HER samples four achieved goals per transition. These values
are fixed across variants unless the named intervention removes a component.

For each policy seed, training proceeds through 300,000 static steps on the
fixed density-0.20 layout, 100,000 steps with one mild dynamic obstacle on the
same-density base layout, and 100,000 steps with the full fixed-toggle
configuration on that same base layout. The environment layout
seed is held at 42 while policy seeds vary, separating optimization variance
from layout variance. Checkpoints, logs, exact configuration, throughput, and
machine metadata are stored for each seed. The policy seeds are 11, 22, 33,
44, and 55. Seven configurations times five seeds yield 35 completed runs and
17.5 million environment interactions.

\subsection{Controlled Ablations}

The complete method is compared with vanilla DQN, HER disabled, potential
shaping disabled, curriculum disabled, full-grid observation, and training on
the full dynamic condition from scratch. Each ablation changes one named
factor. Dynamic-from-scratch training receives 500,000 dynamic steps so its
interaction budget matches the complete curriculum's cumulative budget.

\subsection{Adaptability Metrics}
\label{sec:adaptability}

For every observed change event, we record:
\begin{itemize}
    \item the optimal remaining cost immediately before and after the change;
    \item \emph{optimal-cost shock}, their difference;
    \item \emph{recovery steps}, the decisions until optimal remaining cost
    returns to or below its pre-change value;
    \item event reaction time and node expansions for classical planners, or
    policy decision latency for RL; and
    \item whether the route succeeds after the event.
\end{itemize}
Unrecovered events remain explicitly missing rather than being dropped.
